\documentclass[11pt]{article}
% Shared preamble for the Carry-Adder Wall paper series.
% Each paper does:  % Shared preamble for the Carry-Adder Wall paper series.
% Each paper does:  % Shared preamble for the Carry-Adder Wall paper series.
% Each paper does:  \input{../../shared/preamble}
\usepackage[margin=1in]{geometry}
\usepackage{amsmath, amssymb, amsthm}
\usepackage{microtype}
\usepackage{graphicx}
\usepackage[hidelinks]{hyperref}

\newtheorem{theorem}{Theorem}
\newtheorem{corollary}{Corollary}
\newtheorem{lemma}{Lemma}
\theoremstyle{remark}
\newtheorem*{remark}{Remark}

\newcommand{\Mhash}{\mathsf{M}}
\newcommand{\SHA}{\mathrm{SHA\text{-}256}}
\newcommand{\SHAd}{\mathrm{SHA\text{-}256d}}
\newcommand{\MAJ}{\mathrm{MAJ}}
\newcommand{\Ch}{\mathrm{Ch}}
\newcommand{\Maj}{\mathrm{Maj}}
\newcommand{\bigO}{\mathcal{O}}

\usepackage[margin=1in]{geometry}
\usepackage{amsmath, amssymb, amsthm}
\usepackage{microtype}
\usepackage{graphicx}
\usepackage[hidelinks]{hyperref}

\newtheorem{theorem}{Theorem}
\newtheorem{corollary}{Corollary}
\newtheorem{lemma}{Lemma}
\theoremstyle{remark}
\newtheorem*{remark}{Remark}

\newcommand{\Mhash}{\mathsf{M}}
\newcommand{\SHA}{\mathrm{SHA\text{-}256}}
\newcommand{\SHAd}{\mathrm{SHA\text{-}256d}}
\newcommand{\MAJ}{\mathrm{MAJ}}
\newcommand{\Ch}{\mathrm{Ch}}
\newcommand{\Maj}{\mathrm{Maj}}
\newcommand{\bigO}{\mathcal{O}}

\usepackage[margin=1in]{geometry}
\usepackage{amsmath, amssymb, amsthm}
\usepackage{microtype}
\usepackage{graphicx}
\usepackage[hidelinks]{hyperref}

\newtheorem{theorem}{Theorem}
\newtheorem{corollary}{Corollary}
\newtheorem{lemma}{Lemma}
\theoremstyle{remark}
\newtheorem*{remark}{Remark}

\newcommand{\Mhash}{\mathsf{M}}
\newcommand{\SHA}{\mathrm{SHA\text{-}256}}
\newcommand{\SHAd}{\mathrm{SHA\text{-}256d}}
\newcommand{\MAJ}{\mathrm{MAJ}}
\newcommand{\Ch}{\mathrm{Ch}}
\newcommand{\Maj}{\mathrm{Maj}}
\newcommand{\bigO}{\mathcal{O}}


\title{\textbf{Neural Distinguishers Expire on Carry Composition}}
\author{Peter Hollows}
\date{\today}

\begin{document}
\maketitle

\begin{abstract}
A hand-built carry-aware score measures a one-round advantage cliff on
reduced SHA-256~\cite{paperII}: it predicts a downstream output byte one
adder layer away and loses all reach one round deeper. The natural
objection is that a learned model might find local structure a human
feature misses, as Gohr's neural distinguishers did for round-reduced
Speck~\cite{gohr}. We translate that methodology to the mining read point.
A residual network receives, as features, everything computed through an
interior round (the state words, their carry-free derivations, the round's
modular sums, and the schedule words, in bit, value, and Fourier
encodings) and is free to learn any function of them. The network exceeds
the hand-built score one adder layer downstream ($0.998$ versus $0.88$
retained advantage) and then collapses to the noise floor one round deeper,
across independent stems, a shifted read point, and a $90\times$ capacity,
$64\times$ data, and $8\times$ training-time scaling. On pure $k$-operand
modular sums, with no SHA structure present, the same network learns the
top byte for $k\le3$ and fails for $k\ge4$. We show the wall is not an
artifact of finite feature precision (a seed-paired
\texttt{float32}-versus-\texttt{float64} comparison is null in both arms)
and that it is the learner's reach rather than an intrinsic boundary: the
carry chain's spectral gap is $\tfrac12$ for every $k$~\cite{paperII}, so
nothing in the carry's mixing singles out $k=4$. We also identify a
distinct second failure mode, feature isolation, and show it does not
affect the main result.
\end{abstract}

% =====================================================================
\section{Introduction}
% =====================================================================

The first paper in this series~\cite{paperI} reduces the question of
whether mining beats brute force to a single residual: a global algebraic
shortcut that resolves the output without resolving the carries on the
path. The second~\cite{paperII} measures how far the natural local methods
reach toward such a shortcut, using a hand-built carry-aware score, and
finds a one-round cliff: strong selection of a downstream output byte one
adder layer away, below detection one round deeper, against a $386$-layer
separation.

A single hand-built feature invites one objection above all. A learned
model, searching the local feature space with gradient descent, might find
structure the human score discards. This is precisely what Gohr did to
round-reduced Speck~\cite{gohr}, an ARX cipher with the same
modular-addition nonlinearity: a trained distinguisher beat the
state-of-the-art hand-built differential distinguisher. We run that
experiment at the mining read point and report four findings. The learned
attacker beats the hand score where signal exists and dies at the same
single-round boundary (Sections~\ref{sec:setup}--\ref{sec:scale}). On pure
carry composition, free of any SHA structure, it expires between two and
three chained additions (Section~\ref{sec:ladder}). The wall survives full
double-precision rebuilding (Section~\ref{sec:float}) and is the learner's
reach rather than an intrinsic property of the arithmetic
(Section~\ref{sec:learner}). A separate failure mode, the inability to pull
relevant bits out of distractors, is real but does not drive the headline
(Section~\ref{sec:isolation}).

% =====================================================================
\section{Setup}\label{sec:setup}
% =====================================================================

We use the selection-advantage metric of the anchor sweep~\cite{paperII}:
take the best fraction $1/256$ of candidates by predicted target, and
report $1 - \overline{y}_{\text{selected}}/\overline{y}_{\text{all}}$,
where $y$ is the top byte of the downstream state word. Advantage $0$ is
chance; advantage near $1$ is near-perfect selection. The metric is
identical to the hand-built measurement, so the two attackers are directly
comparable at each prediction depth $j$ (rounds past the read point).

The network is a GELU residual multilayer perceptron of about $7.4$ million
parameters with a $256$-way top-byte head, trained with cross-entropy on
$2^{20}$ candidates and tested on $2^{19}$ at each cell, over $50$ epochs.
At an interior read round it receives, as engineered features, everything
computed through that round: the state words, their carry-free
derivations, the round's modular sums $T_1$ and $T_2$, and the schedule
words, in bit, value, and byte-harmonic Fourier encodings. The protocol was
pre-registered, with amendments recorded as the positive control was
brought to pass.

That positive control is itself informative. The network learns the top
byte of one two-operand modular addition, $\mathrm{topbyte}(T_1+T_2)$ from
Fourier-featured operands, to advantage $0.998$. It failed, at every
capacity, feature, and optimizer setting tried, to compose $T_1$ and $T_2$
from their seven constituent words, a depth of about five adder layers,
until those sums were supplied directly. The learned attacker's
demonstrated arithmetic reach is therefore between one and about five adder
layers, which bounds how the later nulls should be read.

% =====================================================================
\section{The one-round cliff}\label{sec:cliff}
% =====================================================================

Read at an interior round, with the modular sums supplied so the network
starts one adder layer from the target, the learned distinguisher reaches
advantage $0.9978$ (range $0.9972$ to $0.9990$ over three stems), against
the hand-built score's $0.8841$. The network learns the carry-in
corrections that the truncated $(T_1\!\gg\!24)+(T_2\!\gg\!24)$ score
discards, so it is a strictly stronger member of the same local attack
class.

One round deeper, the advantage is gone. Every $j\ge1$ cell sits at the
noise floor. The unbiased final-epoch statistic, pooled over the twelve
$j\ge1$ cells, is $-0.0040$ with $95\%$ confidence interval $[-0.0079,
-0.0001]$; at $j=1$ alone it is $-0.0060$, interval $[-0.0126, +0.0007]$.
A configuration-matched shuffled-label control reads $0.0212$, inside the
$j\ge1$ band, which pins those readings as the pipeline's
maximum-over-epochs noise floor rather than residual signal. The learned
attacker, stronger than the hand score where signal exists, retains no
advantage one round past the read point. The pattern holds at a shifted
read point and across independent stems.

% =====================================================================
\section{Scaling does not move the wall}\label{sec:scale}
% =====================================================================

A null at $j\ge1$ could in principle be a capacity or training limit rather
than a reach limit. We scaled each axis independently at $j=1$, pairing
every cell against a configuration-matched shuffle-label run, with the
verdict being the gap between the real and shuffle maxima (since
maximum-over-epochs advantage inflates with capacity and epoch count, a
fixed gate would manufacture signal as scale grows). Capacity from $0.4$ to
$35.7$ million parameters ($90\times$), data from $2^{16}$ to $2^{22}$
samples ($64\times$), and training time from $50$ to $400$ epochs
($8\times$, a grokking probe~\cite{grokking}) all return null, with no
dose-response on any axis and no late generalization transition. The $j=1$
null is a reach limit: a $90\times$ capacity range, a $64\times$ data
range, and an $8\times$ time range do not move the learned attacker past
one round. (An earlier $500\times$ capacity endpoint was discarded as
invalid, since its test predictions collapsed to a constant from epoch $35$
on; the defensible range is $90\times$.)

% =====================================================================
\section{Where learning stops: the $k$-operand ladder}\label{sec:ladder}
% =====================================================================

The $j\ge1$ nulls confound two readings: that SHA's carry structure
destroys the signal, and that this optimizer cannot compose modular
arithmetic in this format. We separate them on data with no SHA structure
at all. The task is the top byte of a sum of $k$ uniform $32$-bit operands,
with the same pipeline and features. The network solves $k=2$ ($0.9996$)
and $k=3$ ($0.9983$), then fails at $k=4$ and through $k=7$ (maxima
$0.025$ to $0.041$, final-epoch advantages near zero). The learned attacker
expires on carry composition itself, between two and three chained
additions, on synthetic data.

\begin{center}
\begin{tabular}{lcccc}
\hline
operands $k$ & $2$ & $3$ & $4$ & $5$--$7$\\
\hline
max advantage & $0.9996$ & $0.9983$ & $0.041$ & dead ($\le 0.04$)\\
\hline
\end{tabular}
\end{center}

A finer probe sharpens the death point. Holding $k=4$ and sweeping the
operand width, the task is learnable at widths $8$ through $28$ bits and
dead at width $32$, with training bimodal below full width: a seed either
solves the task almost perfectly or sits at the floor, with little in
between, and at width $32$ no seed of many has left the floor. The $k=4$
death is therefore a joint condition, $k\ge4$ and full width, and it has
the character of an optimization cliff whose success probability falls to
zero rather than a smooth capacity limit.

\begin{center}
\begin{tabular}{lccccccc}
\hline
operand width (bits) & $8$ & $12$ & $16$ & $20$ & $24$ & $28$ & $32$\\
\hline
seeds reaching ${\sim}1$ & $1/2$ & $2/2$ & $1/2$ & $1/2$ & $3/3$ & $1/2$ & $\mathbf{0/8}$\\
\hline
\end{tabular}
\end{center}

% =====================================================================
\section{The wall survives double precision}\label{sec:float}
% =====================================================================

The bimodal width result admits a precision reading: the value and Fourier
features resolve an operand to about $2^{-24}$ in single precision, so an
analog shortcut (estimate the real-valued sum, read its top byte) would
lose the carry-relevant low bits exactly when the operand width exceeds the
mantissa, near width $32$. We tested this directly. For each seed we built
the $k=4$, width-$32$ task twice, with identical integer operands and
identical parameter initialization, differing only in numeric precision:
full double precision, features and network arithmetic, against the single-
precision baseline. (Full double precision is required, since any cast back
to single precision after the features reintroduces the same $2^{-24}$
wall.)

The two arms are null together. Across six seed-paired runs, both
precisions leave the advantage at the floor, $0/6$ seeds learning in
either, with the double-precision arm tracking the single-precision arm to
within about $0.01$ at every seed and a maximum advantage of $0.086$.
Restoring full operand resolution does not let gradient descent find the
carry circuit. On this evidence the wall is a trainability barrier, not an
artifact of finite feature precision.

% =====================================================================
\section{The learner's reach}\label{sec:learner}
% =====================================================================

The $k=4$ threshold is a property of the learner, not of the carry chain.
The carry chain for $k$ operands has contraction rate $\tfrac12$ for every
$k$~\cite{paperII}: its eigenvalues are $\{2^{-j}\}$, and larger $k$ adds
only faster-decaying modes, so nothing in the carry's own mixing
distinguishes $k=4$ from any other operand count. The intrinsic
correlation decay is set by the base, not by $k$. What the ladder isolates
is the reach of gradient descent over fixed features at composing carries,
a statement about this attack class.

This is the same conclusion the precision result reaches by a second route.
Together they say the $k=4$ cliff is where the gradient-descent attacker's
last analog shortcut runs out, not an information boundary in the
arithmetic. The carry, the one part of modular addition with no low-degree
or analog surrogate, is what the learner cannot synthesize; this is
consistent with Gohr's distinguishers, which exploited statistical bias
already present in the data they were given rather than synthesizing cipher
arithmetic.

% =====================================================================
\section{A second failure mode: feature isolation}\label{sec:isolation}
% =====================================================================

Probing the ladder surfaced a separate way the learner fails, which we
record so it is not confused with carry depth. Holding $k=4$ and full width
and varying which byte of the sum is predicted, every read-out is dead,
including the low byte. The low byte of the sum is, mathematically, an
$8$-bit modular addition of the operands' low bytes, the same easy task a
width-$8$ cell learns, so its death is not about carry depth. Varying the
feature set isolates the cause: with all features, or with all $32$ bits
and no analog features, the low byte stays dead; with only the eight
relevant low bits, a seed revives to advantage $0.45$. The killer is the
presence of $24$ distractor bits, not the analog features. The learner
cannot extract the relevant sub-circuit when its inputs are buried among
irrelevant ones.

This is a weakness of a multilayer perceptron over a fixed, wide feature
set, in principle addressable by attention or feature selection, so it is
not an intrinsic-hardness claim. It does not touch the main result: the
$k$-operand ladder and the width ladder use top-byte targets, and the top
byte of a $k$-operand sum depends on every operand bit through the carry
chain, so no input bit is a distractor there. A Gohr-faithful
convolutional architecture over bit positions was also run and reached only
$0.03$ to $0.06$ on a task the perceptron passes above $0.99$; its
inductive bias is matched to the \textsc{xor}-differential locality Gohr's
distinguishers exploit, which carry chains do not have.

% =====================================================================
\section{Scope, reproduction, and conclusion}\label{sec:scope}
% =====================================================================

This tests a one-shot learned predictor, the direct analogue of Gohr's
distinguisher. It does not rule out machine learning used as a search
heuristic over algebraic representations; Gohr's strongest attacks combined
the network with classical key-ranking and multi-round search, which is the
same residual the first paper~\cite{paperI} scopes out and does not claim
to close. The experiment uses the simplified single-compression read point
of the anchor sweep~\cite{paperII}, faithful to the round-function
mechanism but not the full mining map, and ``null'' means below a
pre-registered gate at the tested power, consistent with zero rather than a
proof of zero.

The experiment is reproducible from a pre-registered protocol with
documented amendments, training code, and raw per-cell results; unlike the
two static results of~\cite{paperII} it requires a GPU and the JAX/optax
stack. The vendored $\SHA$ is verified bit-exact against \texttt{hashlib}.

A learned distinguisher of the kind that broke round-reduced Speck beats
the hand-built carry-aware score one adder layer downstream and retains no
advantage one round deeper. It expires on carry composition
between two and three chained additions, a wall that is independent of
feature precision and is the reach of the learner rather than a boundary in
the arithmetic. The one-round cliff of~\cite{paperII} is not an artifact of
a weak hand-built score, and the $386$-layer separation stands against the
strongest local attack we could train.

% =====================================================================
\begin{thebibliography}{99}

\bibitem{paperI}
P.~Hollows.
\emph{The Amortization Envelope: Where a SHA-256d Mining Advantage Can and
Cannot Live.}
Carry-Adder Wall series~I, 2026.

\bibitem{paperII}
P.~Hollows.
\emph{Carry Depth: A Circuit-Independent Coordinate of ARX, and the
One-Round Advantage Cliff.}
Carry-Adder Wall series~II, 2026.

\bibitem{gohr}
A.~Gohr.
\emph{Improving Attacks on Round-Reduced Speck32/64 Using Deep Learning.}
CRYPTO 2019.

\bibitem{grokking}
A.~Power, Y.~Burda, H.~Edwards, I.~Babuschkin, V.~Misra.
\emph{Grokking: Generalization Beyond Overfitting on Small Algorithmic
Datasets.}
2022, \texttt{arXiv:2201.02177}.

\end{thebibliography}

\end{document}
